% §8 Conclusion — Target 0.5pp

\section{Conclusion}
\label{sec:conclusion}

We presented the first empirical comparison of four Bitcoin
covenant proposals for vault custody---\opctv{}, \opccv{},
\opvault{}, \catcsfs{}---plus a Simplicity vault on Elements
(Appendix~\ref{app:simplicity}) as a cross-substrate reference that
occupies the dominant corner. Two structural results constrain the
design space. Proposition~\ref{thm:griefing_safety} shows that no
covenant vault simultaneously achieves permissionless recovery and
griefing resistance. The design-space decomposition of
Section~\ref{sec:bg:axes} shows that vault designs factor along
seven canonical axes, with six derived attack classes
\classfee{}--\classmode{} each a deterministic consequence of a
specific axis-value (Sections~\ref{sec:imposs:class_fee}--\ref{sec:imposs:class_mode}).
As a structural consequence, no proposed Bitcoin extension
sits at the corner of the space where defender loss is
simultaneously zero against all six classes; each falls short on
at least one axis with a corresponding class surface. The
\axauth{}--\axintro{} axes provide a lookup procedure for localising
future BIPs without additional measurement.

The cost of custody is therefore not a single number. It depends on
the covenant's axis tuple, the class the operator prioritises, the
prevailing fee regime, and the watchtower's operational strategy.
Under two designs equal at the script layer, the relative safety can
still invert with the defender's strategy: safety rankings among
current covenants are not operator-free. No proposal minimises the
cost of custody against every class; the seven-axis decomposition
makes the tradeoffs precise and identifies a structurally-safer
corner $D^\star$ that the Bitcoin BIP corpus has not yet occupied.
Open directions include mechanised proofs of the class derivations,
the impact of TRUC/BIP-431 and cluster-mempool on class \classfee{},
temporal verification of class \classhot{} races, and cross-chain
generalisation of the decomposition. The framework, all measurements,
formal models, and both original vault implementations are available
at the anonymised artifact (URL withheld for double-blind review).

\providecommand{\todo}[2][]{\textbf{[TODO: #2]}}
